\chapter{课堂之中}\label{s:classroom}

上一章讲了如何练习授课，
并介绍了一种方法——现场编程——
让老师能适应学习者的节奏和兴趣。
本章介绍其他一些在编程课上同样有用的做法。

在介绍它们之前，
值得先停一下，设定好预期。
我们知道的最好的教学方法是个别辅导：\index{个别辅导（的效果）}
\cite{Bloo1984}发现，接受一对一教学的学生，
比接受传统讲授的学生高出两个标准差，
也就是接受个别辅导的学生胜过
98\% 接受讲授的学生。
然而，
尽管在历史上大部分时期，师徒制和学徒制是最常见的知识传授方式，
正规教育的工业化却使它在今天成了例外。
尽管人工智能被吹得天花乱坠，\index{人工智能（炒作）}
它短期内也不会把这个矛盾化解掉，
所以下面描述的每一种方法，本质上
都是想在规模化的情况下逼近个别辅导的效果。

\seclbl{执行行为准则}{s:classroom-coc}

过去三十年里，我在教学上最重要的一条体会是，
不管课内课外，每个人都尊重其他人，是多么重要。
如果你以任何方式使用这些材料，
请采用一份像\appref{s:conduct}里那样的行为准则，\index{行为准则}
并要求所有参与你课程的人都遵守它。
它拦不住人出言伤人，
就像反盗窃的法律拦不住人偷东西一样，
但它\emph{能}把预期和后果讲清楚，
并表明你在努力让你的课堂对所有人都友善。

但行为准则只有被执行才有用。
如果你认为有人违反了你的准则，
你可以根据违规的严重程度、以及你是否认为它是有意的，
警告他们、
要求他们道歉、
和/或将他们请出课堂。
不管你做什么：\index{行为准则!执行}

\begin{description}

\item[在有证人在场的情况下做。]
  大多数人在有听众时都会收敛言辞和敌意，
  而且有别人在场，能确保
  事后的讨论不会退化成「谁说了什么」的互相矛盾的口角。

\item[如果你把某人请出课堂，就向全班说明并解释原因。]
  这有助于防止谣言传播，
  并表明你的行为准则确实是有分量的。

\item[尽快给违规者发邮件，]
  概括发生了什么以及你采取了哪些步骤，
  并把邮件抄送给你的工作坊主办方或某位同行老师，
  这样就能留下当时的对话记录。
  如果违规者回复了，
  不要陷入冗长的争辩：
  那从来不会有成果。

\end{description}

课外发生的事，至少和课内发生的事一样重要~\cite{Part2011}，
所以你需要提供一种途径，让学习者能报告你自己看不到的问题。
一个办法是请一位不属于你们小组的人做第一联系人；
这样，
如果有人想投诉你或你的某位同行，
他们就对保密性和独立处理有了一定保障。
\cite{Auro2019}还有很多别的建议，
既简短又实用。

\seclbl{同伴教学}{s:classroom-peer}

不管一位老师有多好，
他们一次也只能讲一件事。
那么他们怎么才能在合理的时间内清除许多不同的误解呢？
到目前为止发展出来的最好办法，是一种叫\gref{g:peer-instruction}{同伴教学}的技巧。
它最初由哈佛的 Eric Mazur 提出~\cite{Mazu1996}，\index{Mazur, Eric}
在各种情境下被广泛研究，
包括编程~\cite{Crou2001,Port2013}，
而~\cite{Port2016}发现，即便第一次接触，学习者也很看重同伴教学。

同伴教学试图以一种可扩展的方式提供一对一教学，
办法是把形成性评价和学习者讨论交错起来：

\begin{enumerate}

\item
  对这个主题做简要介绍。

\item
  给学习者一道多项选择题，用来探查他们的误解
  （而不是测试简单的事实性知识）。

\item
  让所有学习者就 MCQ 的答案投票。

  \begin{itemize}

  \item
    如果学习者全都答对了，就往下走。

  \item
    如果他们全选了同一个错误答案，
    就针对那个具体误解加以讲解。

  \item
    如果对错混杂，
    就给他们几分钟，让他们 2—4 人一组互相争辩，
    然后再投一次。

  \end{itemize}

\end{enumerate}

正如
\hreffoot{https://www.youtube.com/watch?v=2LbuoxAy56o}{这段视频}所示，
小组讨论会显著改善学习者的理解，
因为它暴露出他们推理中的空缺，并迫使他们把自己的想法讲清楚。
之后重新投票，就能让老师知道是否可以往下走，
还是还需要进一步解释。
在给出正确答案之后再补一轮额外讲解，
又给了学习者一次巩固理解的机会。

但这会不会是假象？
成绩提升，是因为讨论中理解加深了，
还是仅仅因为「跟着领头的走」（「跟着 Jane 投，她总是对的」）？
\cite{Smit2009}通过在第一道题之后紧跟一道学习者独立作答的题来检验这一点。
他们发现，同伴讨论确实能加深理解，
即便讨论小组里原本没有一个人知道正确答案。
只要组内存在意见分歧，
他们的误解就会互相抵消。

\newpage
\begin{aside}{公开站定立场}
  让学习者公开投票很重要，
  这样他们就不能事后改主意、并靠自我开脱
  （比如「我刚才只是看错题了」）来给自己找理由。
  同伴教学很大一部分价值，来自答错之后的超校正、
  以及不得不把原因想透的过程
  （\secref{s:individual-strategies}）。
\end{aside}

\seclbl{协同教学}{s:classroom-together}

\grefdex{g:co-teaching}{协同教学}{协同教学}指两位老师在同一间教室里一起工作的任何情形。
\cite{Frie2016}描述了几种做法：

\begin{description}

\item[双师教学：]
  两位老师协同讲授同一份内容流，
  像音乐家轮流独奏那样交替进行。

\item[主讲加辅助：]
  A 老师主讲，B 老师在教室里走动，
  帮助遇到困难的学习者。

\item[分层教学：]
  A 老师给一小群学习者做更密集或更专门的讲授，
  而 B 老师给大组上常规课。

\item[主讲加观察：]
  A 老师主讲，B 老师观察学习者，
  收集他们理解情况的资料，用来帮助规划后续课程。

\item[并行教学：]
  把全班分成两组，
  两位老师同时向各自那组呈现同样的材料。

\item[站点教学：]
  把学习者分成小组，
  在各个站点或活动之间轮转，
  老师在需要的地方进行督导。

\end{description}

所有这些模式，都比独自教学创造了更多非预期知识传递的机会。
\index{非预期知识传递}
双师教学在历时一整天的工作坊里尤其有好处：
它让每位老师的嗓子有机会休息，
也降低了他们一天下来累到开始冲学习者发火、
或在键盘上手忙脚乱的风险。

\begin{aside}{帮忙的人}
  很多不太适合讲课的人，
  却愿意、也能够提供课堂上的技术支持。
  他们可以帮学习者做安装和配置、
  在练习中回答技术问题、
  巡视教室以发现可能需要帮助的人、
  或盯着共享笔记（\secref{s:classroom-notetaking}），
  并回答问题，
  或提醒老师在休息时来回答。

  助手有时是正在受训成为老师的人
  （也就是「主讲加辅助」模式里的 B 老师），
  但他们也可以是主办院校的技术支持人员、
  课程的往届学员、
  或者已经很熟悉材料的进阶学习者。
  让后一种人当助手有双倍效果：
  他们不仅更可能理解同伴遇到的困难，
  也能让这些进阶学习者不至于无聊。
  这有助于全班保持投入，因为无聊是会传染的：
  如果有一小撮人开始走神，
  他们周围的人也会跟着走神。
\end{aside}

如果你和一位搭档协同教学：

\begin{itemize}

\item
  每次上课前花 2—3 分钟
  确认谁讲什么。
  如果有时间，
  可以一起画或过一遍一张概念图。

\item
  也利用这段时间商定几个手势。
  「你讲太快了」、
  「大声点」、
  「那个学习者需要帮助」、
  以及「该上厕所了」都很有用。

\item
  每人一次至少讲 10—15 分钟，
  因为切换太频繁会分散学习者注意力。

\item
  没在讲的那个人不应该打断、
  插话纠正或补充、
  或做任何分散主讲者注意力的事。
  唯一的例外是，如果学习者看起来无精打采或缺乏自信，可以问一些引导性的问题。

\item
  每个人在开始讲课之前，应该花几分钟看看搭档讲完之后要讲什么，
  然后\emph{不要}讲那部分材料。

\item
  没在讲的那个人应该保持对课堂的投入，
  而不是去补看邮件。
  盯着共享笔记（\secref{s:classroom-notetaking}）、
  留意哪些学习者在挣扎、
  记下一些反馈留到下次休息时给你的搭档——任何
  对课堂有贡献的事，都比没有贡献的事强。

\end{itemize}

最重要的是，
下课后花几分钟互相祝贺或互相安慰：
教学一如人生，
痛苦分担忧减，快乐分享增。

\seclbl{评估先备知识}{s:classroom-prior}

你在开课前对学习者了解得越多，
就越能帮到他们。
在正规学校体系里，
你课程的前置要求会告诉你，
他们可能已经知道什么。
但在散养式环境里，
你的学习者可能差异大得多，
所以你可能想在开课前给他们一份简短的调查或问卷，
弄清他们有哪些知识和技能。

让人按 1 到 5 给自己打分是没有意义的，
因为一个人对某个学科了解得越少，\index{自我评估（的危险）}
就越无法准确估计自己的知识
（\figref{f:classroom-dunning-kruger}，
出自 \hreffoot{Neurologica}{https://theness.com/neurologicablog/index.php/misunderstanding-dunning-kruger/}），
这一现象叫做\gref{g:dunning-kruger-effect}{邓宁—克鲁格效应}~\cite{Krug1999}。
反过来，
来自代表性不足群体的人，往往会低估自己的技能。

\figimg{figures/dunning-kruger.png}{邓宁—克鲁格效应}{f:classroom-dunning-kruger}{}

与其让人自我评估，
不如问他们完成某些具体任务能有多容易。
但这么做有风险，
因为学校把人训练成了
把任何看起来像考试的东西都当成必须通过的关卡，
而不是塑造教学的机会。
如果有人在你的前测里哪怕只有那么几道题回答「我不知道」，
他们可能就断定你的课对他们来说太深了。
这样你就可能把很多你最想帮的人吓跑。

\secref{s:checklists-preassess}给出了一份简短的课前问卷，
大多数潜在学习者不容易被它吓到。
如果你用它或类似的东西，
尽量去跟进那些\emph{没有}作答的人，弄清原因，
并把你对学习者的评估和他们的自我评估作比较，
从而改进你的问题。

\seclbl{为能力参差做规划}{s:classroom-mixed}
\index{能力参差（如何应对）}

如果你的学习者先备知识水平相差很大，
你很容易陷入一种局面：班上有三分之一的人跟不上，
有三分之一的人觉得无聊。
这让谁都不满意，
但你可以用一些策略来应对这种情况：

\begin{itemize}

\item
  办工作坊之前，
  通过展示几道他们会做的练习，把难度水平清楚地告诉大家。
  这比列成要点的主题清单更能有效帮潜在参与者判断课程的难度。

\item
  提供额外的自主进度练习，
  让进阶学习者不会早早做完、然后觉得无聊。

\item
  留意掉队的学习者，
  尽早干预，免得他们受挫放弃。

\item
  请进阶学习者帮助旁边的人
  （见下文\secref{s:classroom-pair}）。

\end{itemize}

应对能力参差的另一种办法，
是让每个人像上在线课那样、按自己的节奏学材料，
但大家是同时、并排地进行，
助手在教室里走动，帮人解卡。
以这种方式办工作坊时，有些人会比其他人多走三四倍的进度，
但每个人都会度过有收获、有挑战的一天。

\begin{aside}{假性新手}
  \gref{g:false-beginner}{假性新手}指
  以前学过一门语言、现在又重新学的人。
  在课前测试里，他们可能和\grefdex{g:absolute-beginner}{绝对新手}{绝对新手}
  难以区分，
  但一旦开课，他们就能快得多地推进，
  因为他们是在重新学，而不是第一次学。

  是假性新手，常常是\gref{g:preparatory-privilege}{预备性特权}的一个标志~\cite{Marg2010}，
  假性新手在散养式编程班里很常见。
  例如，
  一个家境富裕到能送他去机器人夏令营的孩子，
  在编程知识前测里可能成绩很差，
  因为那些材料在他脑子里已经不新鲜了，
  但他仍然比家境没那么好的孩子更有优势。
  上面描述的策略能在这种情况下有助于拉平起点，
  但再说一遍，
  真正的解决办法是用你自己的特权
  去处理课堂之外、更大的因素~\cite{Part2011}。
\end{aside}

最重要的一点是接受：
你没法时时刻刻帮到每个人。
如果你为了照顾两个跟不上的学生而慢下来，
你就辜负了另外十八个。
同样，
如果你花几分钟和一个觉得无聊的学习者讲进阶主题，
班里其他人就会觉得被冷落。

\seclbl{结对编程}{s:classroom-pair}

\grefdex{g:pair-programming}{结对编程}{结对编程}是一种软件开发实践，
两个人\hreffoot{https://www.youtube.com/watch?v=vgkahOzFH2Q}{一起在一台电脑上工作}。
一个人（驾驶员）负责打字，
另一个人（领航员）提出意见和想法，
两人每小时交换几次角色。

结对编程是专业工作中一种有效的实践~\cite{Hann2009}，
也是一种很好的教学方式：
它的好处包括入门课程通过率更高、
软件质量更好、
学习者对自己解决方案更有信心。
还有证据表明，来自代表性不足群体的学习者
受益甚至更多~\cite{McDo2006,Hank2011,Port2013,Cele2018}。
搭档可以在实操练习中互相帮助、
在讲解答案时澄清彼此的误解、
在休息时讨论共同的兴趣。
我发现它在能力参差的班里特别有用，
因为结对的人比单独的个体更同质。

使用结对时，
把\emph{每个人}都编入结对，
而不只是那些跟不上的学习者，
这样没人会觉得被单独挑出来。
定期让人换座位
（从而和不同的搭档结对）也很有用，
并让每对中的人每小时交换三四次角色，
免得每对里个性更强的那个人主导整场练习。

如果你的学习者第一次接触结对编程，
花几分钟演示一下它实际是什么样子，
好让他们明白
手不在键盘上的那个人
不该只是坐着看。
最后，
告诉他们，那些一心想着尽快完成任务的人，
在分享上会更不公平~\cite{Lewi2015}。

\begin{aside}{交换搭档}
  对于是否应该要求人们定期更换搭档，老师们意见不一。
  一方面，它给了每个人获得新洞见、结交新朋友的机会。
  另一方面，
  一天里好几次把电脑和电源适配器搬到新桌子上会造成干扰，
  而且结对对内向的人可能不太自在。
  话虽如此，
  \cite{Hann2010}发现「大五」人格特质
  和结对编程中的表现之间只有微弱相关，
  尽管一项更早的研究~\cite{Wall2009}发现，
  成员在人格特质上水平不同的结对，交流得更频繁。
\end{aside}

\seclbl{一起\ldots 记笔记？}{s:classroom-notetaking}
\index{记笔记}

记笔记是一种实时的精细加工（\secref{s:individual-strategies}）：
它迫使你在材料进来时就加以组织和反思，
这反过来又增加了你把它转入长期记忆的可能性。
许多研究表明，
学习时记笔记能改善保持~\cite{Aike1975,Boha2011}。
虽然还没有被广泛研究~\cite{Ornd2015,Yang2015}，
但我发现，让学习者在一个共享的在线页面上一起记笔记也同样有效：

\begin{itemize}

\item
  它让人能把自己以为听到的东西
  和别人听到的东西作比较，
  这有助于他们当场填补空缺、纠正误解。

\item
  它给了班里进阶的学习者一些有用的事做。
  与其在课上觉得无聊去刷 Instagram，
  他们可以带头记录正在讲的内容，
  这让他们保持投入，
  也让不太进阶的学习者能把更多注意力放在新材料上。

\item
  学习者记的笔记，通常比老师提前准备好的笔记\emph{对他们}更有帮助，
  因为学习者更可能写下他们真正觉得新的东西，
  而不是老师预测会显得新的东西。

\item
  在学习者做练习时瞄一眼最近的笔记，
  能帮老师发现全班漏掉了或误解了什么。

\end{itemize}

\begin{aside}{笔真比键盘更强吗？}
  \cite{Muel2014}报告说，用电脑记笔记
  通常不如用纸笔记笔记有效。
  尽管他们的结果被广泛传播，
  \cite{More2019}却没能重复出来。
\end{aside}

如果学习者在一起记笔记，
你也可以让他们把代码小片段
和对形成性评价问题的要点式或句子式回答粘贴进去。
为了防止所有人同时去编辑那几行，
把每个人的名字列成一份清单，
在你想让每个人都回答一个问题时，把清单粘贴到文档里。

学习者第一次尝试一起记笔记时，常会觉得它让人分心，
因为他们得把注意力分配到
老师正在说的
和同伴正在写的之间（\secref{s:architecture-brain}）。
如果你只是和某个特定群体合作一次，
那么就应该听从\secref{s:classroom-innovate}里的建议，
让他们各自单独记笔记。

\begin{aside}{可改进之处}
  向学习者表明他们是在和你\emph{一起}学、
  而不只是\emph{从}你这里学，一个办法
  是允许他们通过编辑你的课程（副本）来记笔记。
  与其发布 PDF 让他们下载，
  不如把你的幻灯片、笔记和练习的可编辑副本
  放到 wiki、
  Google Doc，
  或任何允许你审阅并评论改动的地方。
  给指出错误、
  讲清解释、
  添加新例子、
  和编写新练习的人记上一笔功劳，并不会减少你的工作量，
  却能提高参与度和课程的寿命
  （\secref{s:process-maintainability}）。
\end{aside}

\seclbl{便利贴}{s:classroom-sticky-notes}

便利贴是我最喜欢的教学工具之一，
而且不只我一个人喜欢它那多用途、便携、
好贴、好折、
还有那微妙而诱人的气味~\cite{Ward2015}。

\subsection*{作为状态旗}
\index{便利贴!作为状态旗}

给每个学习者两张不同颜色的便利贴，
比如橙色和绿色。
它们可以用来举起来投票，
但真正的用途是当状态旗。
如果有人做完了练习、想让人检查，
就把绿色便利贴贴在笔记本电脑上；
如果遇到问题、需要帮助，
就举起橙色的那张。
这比让人举手好得多：
它更不显眼（这意味着他们更可能真的这么做），
他们的旗子举着时还能继续干活，而不是一只手打字，
老师从教室前面就能一眼看出全班处于什么状态。
当能力参差班里的人按自己的进度学材料时
（\secref{s:classroom-mixed}），状态旗尤其有用。

等你的学习者习惯了两种贴纸，
再给他们第三种，让他们在脑子满了
或需要上厕所时举起来\footnote{一位同事曾告诉我，
教学的基本单位是膀胱。
我说我从没这么想过，
她说：「你显然没怀过孕。」}。
再说一遍，
成年人贴一张便利贴比举手更乐意，
而且一旦有一张蓝色便利贴举起来，
通常就会跟着冒出一片。

\subsection*{用来分配注意力}
\index{便利贴!用来分配注意力}

便利贴也可以用来确保老师的注意力得到公平分配。
让每个学习者把名字写在便利贴上，
贴在自己的笔记本上。
老师每叫到他们一次、或回答了他们的一个问题，
他们就把自己的便利贴取下来。
等所有便利贴都取下来之后，
大家再把自己的重新贴上去。

这个技巧让老师很容易看出自己最近没和谁说过话，
进而帮助他们避免无意识的偏见、
避免优先和最外向的学习者互动。
没有这样的检查，
就太容易形成一种反馈循环：外向的人得到更多关注，
于是表现更好，
于是又得到更多关注，
而更安静、更不自信或边缘化的学习者则被落在后面~\cite{Alvi1999,Juss2005}。

它也向学习者表明，注意力在被公平地分配，
这样他们\emph{被}叫到时，
就不会觉得自己是被针对。
当我和一个新群体合作时，
在课上的头一两个小时里，我会允许人们在自己不愿被叫到时
取走自己的便利贴。
如果他们随着时间推移一直这么做，
我会试着和他们私下聊一聊，弄清原因，
并看看我能不能做点什么让他们更自在。

\subsection*{作为一分钟卡片}
\index{便利贴!作为一分钟卡片}

你也可以把便利贴当作\grefdex{g:minute-cards}{一分钟卡片}{一分钟卡片}。
每次休息前，
学习者花一分钟，在绿色便利贴上写下一件他们觉得有用的东西，
在橙色便利贴上写下一件他们觉得太快、
太慢、
太乱、
或无关的东西。
他们去喝咖啡或吃午饭时，
你翻看这些纸条，找找规律。
在一个 40 人的班里，不到五分钟就能看出学习者在享受什么、
被什么搞糊涂了、
遇到什么问题、
以及你还有哪些问题没回答。

学习者不应该在一分钟卡片上签名：
它们本意是匿名反馈。
\secref{s:classroom-practices}里描述的「一条上、一条下」技巧，
则是获得集体、可署名反馈的机会。

\seclbl{永远不要留空白页}{s:classroom-blank}

编程工作坊和其他种类的课程，
可以围绕一组相互独立的练习来构建，
也可以分阶段展开一个单独的扩展例子，
或者采用混合策略。
独立练习的两大优点是：
掉队的人能很容易重新跟上，
课程开发者可以随意增删和重排材料
（\secref{s:process-maintainability}）。
而一个单独的扩展例子，
则能向学习者展示他们正在学的零碎东西如何拼合在一起：
用教育学的说法，
它为他们整合知识提供了更多机会。

不管你采用哪种做法，
新手都不应该对着一张空白纸或空白屏幕开始做练习，
因为他们常会觉得这令人发怵或不知所措。
如果他们一路跟着你做了现场编程，
就让他们再多加几行，
或者改动你搭建起来的例子。
或者，如果他们在一起记笔记，
就往文档里粘贴几行起始代码，让他们去扩展或修改。

修改现有代码，而不是从头写新代码，
不仅给了学习者结构：
它也更接近他们在真实生活中会做的事。
不过要记住，
学习者可能会因为试图弄懂所有起始代码而分心，
反而没做自己的练习。
Java 的 \texttt{public static void main()}
或 Python 程序顶部那几行 \texttt{import} 语句，
对你来说可能讲得通，
但对他们来说是多余的负荷（\chapref{s:architecture}）。

\seclbl{为学习者搭好环境}{s:classroom-setup}

散养式学习者常常想自带电脑，
并且希望离开课堂时，这些电脑已经配置好、能干真活。
所以散养式老师应该准备好同时在 Windows 和 MacOS 上教学\footnote{「还有 Linux！」
有人从教室后面喊道。}，
尽管只要求学习者用一种系统会简单得多。

\begin{aside}{共同点}
  如果你的参与者用的是不同的操作系统，
  就尽量避免使用只属于某一种系统的功能，
  并把你\emph{确实}用到的这类功能指出来。
  例如，
  Windows 上「最小化窗口」的控件和行为
  就和 MacOS 上不一样。
\end{aside}

不管你要应付多少种平台，
都要把详细的配置说明放到课程网站上，
并在工作坊开始前几天给学习者发邮件，
提醒他们提前配置好。
仍然会有少数人没装好所需软件就来，
因为他们遇到了问题、
抽不出时间完成所有步骤、
或者就是那种从不提前照说明做的人。
为了发现这种情况，
让大家一到就运行某条简单命令并把结果给老师看，
然后让助手和其他学习者
去帮助遇到麻烦的人。

\begin{aside}{虚拟机}
  有些人用 \hreffoot{http://docker.com}{Docker} 之类的工具
  在学习者电脑上装虚拟机，\index{虚拟机}
  好让每个人用的工具完全一样，
  但这会引入一堆新问题。
  又老又小的机器根本跑不动它们，
  学习者还要在复制粘贴等操作的两套不同键盘快捷键之间来回切换，
  很吃力，
  甚至连胜任者都会搞不清到底哪个操作发生在哪里。
\end{aside}

配置这么麻烦，
很多老师更喜欢让学习者用基于浏览器的工具。
然而，
这会让课程依赖院校的 WiFi，\index{院校 WiFi（的隐患）}
（其质量可能非常参差），
也无法满足学习者「带着自己配置好、能用于真实场景的电脑离开」的愿望。
不过，随着像 \hreffoot{https://glitch.com/}{Glitch}
和 \hreffoot{http://rstudio.cloud}{RStudio Cloud} 这样的云端工具变得越来越健壮，
后一种顾虑也就没那么重要了。

解决配置问题的最后一种办法，是把课程拆到好几天里，
让人在每天结束离开之前，
把第二天需要的东西装好。
把工作分成小块，每一块就没那么吓人，
学习者更可能真的去做，
而且它能保证除了第一天之外，每节课你都能按时开始。

\seclbl{其他教学实践}{s:classroom-practices}

下面描述的这些小做法都不是必需的，
但全都会改善授课。
和国际象棋、婚姻一样，
教学上的成功往往是缓慢而稳定的进步的结果。

\subsection*{从自我介绍开始}

上课时先介绍你自己。
如果你是专家，
就给他们讲讲你是怎么走到今天这一步的；
如果你只比他们领先两步，
就强调你和他们的共同点。
不管你说什么，
你的目标都是让自己更平易近人，
并让他们更相信自己能成功。

学习者也应该互相介绍。
在一打人的班里，
他们可以口头介绍；
在大一些的班里，或者彼此都是陌生人时，
我发现让他们每人往共享笔记里写一两行关于自己的话更好
（\secref{s:classroom-notetaking}）。

\subsection*{配置好你自己的环境}

配置你自己的环境至少和配置学习者的环境一样重要，
但更复杂。
除了有网络访问权限和你将要用的所有软件，
你还应该有一杯水
或一杯茶、咖啡。
这有助于保持嗓子湿润，
但它真正的用途，是在有人问了个难题、
或你忘了接下来要说什么时，给你一个停顿几秒、想一想的借口。
你大概还会想要几支白板笔
和\secref{s:checklists-events}里描述的其他一些东西。

让你日常工作不干扰教学的一个办法，
是在电脑上另建一个账户专门用来教学。
这个第二账户里所有东西都用系统默认，
配上更大的字体和空白背景，
并关掉通知，这样你的教学就不会被弹窗打断。

\subsection*{整天的课程不要留作业}

学了一整天编程的学习者会很累。
如果你再给他们留课后作业，
第二天他们开始上课时也会很累，
所以别留。

\subsection*{不要碰学习者的键盘}

替学习者把东西修好往往很诱人，
但即便你把每一步都讲出来，
这也可能打击他们的积极性，
因为它凸显了他们和你的知识差距。
相反，
把手从键盘上拿开，用言语引导学习者完成他们该做的事：
这会花更长时间，
但更可能真正被他们掌握。

\subsection*{把问题重复一遍}

每当有人在课上提问，
在回答之前先把问题重复给他听，
以确认你听懂了，
也给可能没听清的人一个机会。
当讲课正在被录像或直播时，这一点尤其重要，
因为你的麦克风通常收不到别人说的话。
把问题重复一遍，也给了你一个机会，
把问题引向你更擅长回答的方向{\ldots}

\subsection*{一条上、一条下}

一分钟卡片的一个补充，是在每天结束时征求总结性反馈。
学习者轮流给出关于这一天的要么一条优点、要么一条缺点，
且不能重复已经说过的东西。
禁止重复迫使人说出他们本来可能不会说的话：
等所有「安全」的反馈都说完了，
参与者就会开始说出他们真正的想法。

\begin{aside}{不同方式，不同回答}
  一分钟卡片（\secref{s:classroom-sticky-notes}）是匿名的；
  交替的上、下反馈则不是。
  你应该把两者结合使用，
  因为匿名既能带来诚实，也可能带来捣乱。
\end{aside}

\subsection*{让学习者做预测}

研究表明，如果被要求预测接下来会发生什么，
人们能从演示中学到更多~\cite{Mill2013}。
这很自然地融入现场编程：
在往程序里加了几行、或改了几行之后，
问全班它运行起来会发生什么。
如果例子哪怕只是中等复杂，
预测也可以作为一轮同伴教学里引人入胜的问题。

\subsection*{摆放桌子}

你可能没法控制上课房间里桌椅的布局，
但如果你能，
我们发现最好是平摆（就餐式）座位，
而不是阶梯（剧场式）座位，
这样你更容易够到需要帮助的学习者，
学习者也更容易互相结对（\secref{s:classroom-mixed}）。
地下电源插座能让你不必把电源线拉过地板，
既让生活更方便，也更安全，
但仍然不常见。

不管你是什么布局，
都要尽量保证每个座位都能无遮挡地看到屏幕。
良好的靠背支撑也很重要，
因为人们会在椅子上坐很久。
和地下电源插座一样，
好的教室座椅至今仍然不幸地不常见。

\subsection*{润喉糖}

如果你一整天对着一屋子人讲话，
嗓子会变哑，因为你在刺激喉部和咽部的上皮细胞。
这不仅让你声音沙哑——还会让你更容易感染
（这也是人们教完课后常常感冒的原因之一）。

保护自己免受此害的最好办法，是让嗓子保持润滑，
而最好的做法是早点、经常地含润喉糖。
好的润喉糖还能掩盖咖啡口气的出现，
为此你的学习者大概会很感激。

\subsection*{思考—结对—分享}

\grefdex{g:think-pair-share}{思考—结对—分享}{思考—结对—分享}是一种轻量的技巧，
帮人们通过与同伴讨论来改进想法。
每个人先独立思考一个问题、并记几笔笔记。
然后两人一组互相解释自己的想法，
把它们合并起来，或挑出最有希望的那个。
最后，
几组把他们的想法展示给全班。

思考—结对—分享之所以有效，是因为它迫使人们把认知外化
（\secref{s:memory-concept-maps}）。
它还让他们有机会在\emph{向}更大的群体公开\emph{之前}，
发现自己想法里的空缺或矛盾，
这能让不太外向的学习者没那么怕出丑。

\subsection*{早、中、晚}

\cite{Smar2018}发现，
如果学习者的课程和其他任务被安排在不吻合他们自然生物钟的时间，
他们表现会更差，
也就是如果「早起型」的人上晚课、或反过来，
他们的成绩会受损。
在小班里通常没法照顾到这一点，
但大一些的班应该试着给并行的场次错开开始时间。
这也能帮到那些要兼顾育儿责任和其他限制的人，
并缩短咖啡休息和上厕所时排队的长度。

\subsection*{幽默}

讲课时应该少用幽默：
大多数笑话写下来就没那么好笑了，
而且每重读一遍就更不好笑。
讲课时即兴的幽默通常效果更好，
但也很容易出岔子：
在你朋友圈里是个笑话的东西，
到你的听众那儿可能变成严肃的政治议题。
如果你讲课时确实要开玩笑，
不要拿任何群体、
或任何个人开涮，也许除了你自己。

\seclbl{限制创新}{s:classroom-innovate}

本章介绍的每一种技巧都会让你的课更好，
但你不应该想一次全盘采用。
原因是每一种新做法都会增加\emph{你}的认知负荷，也会增加学习者的，
因为你们现在既要学一种新的学习方式，
又要学课程的主题内容。
如果你是反复和同一个群体合作，
可以每隔几节课引入一种新技巧；
如果你只是和他们办一天的工作坊，
最好只挑一种他们没见过的做法，
让他们先适应它。

\seclbl{练习}{s:classroom-exercises}

\exercise{制作一份问卷}{个人}{20}

以\secref{s:checklists-preassess}里的问卷为模板，
制作一份你可以在自己上课前发给学习者的简短问卷。
你最想了解他们的什么背景？
双方怎么才能确保，对于你在问的理解水平，大家有共识？

\exercise{你自己的一个做法}{全班}{15}

想一个到目前为止还没被介绍过的教学做法。
把你的想法讲给同伴听，
也听他的，
然后挑一个展示给全班。
（这个练习就是思考—结对—分享的一例。）

\exercise{能让我开吗？}{两人一组}{10}

和同伴交换电脑
（最好是和你用不同操作系统的人），
过一遍一道简单的编程练习。
这有多让人抓狂？
它让你对新手一直要经历的处境多了多少理解？

\exercise{结对}{两人一组}{15}

看\hreffoot{https://www.youtube.com/watch?v=vgkahOzFH2Q}{这段结对编程视频}，
然后和同伴一起练习。
记得每隔几分钟就在驾驶员和领航员之间交换角色。
你们要多久才能进入一个顺畅的节奏？

\exercise{比较笔记}{小组}{15}

分成 3—4 人的小组，
比较你们在本章做的笔记。
你认为值得一提、而同伴漏掉的是什么？反过来呢？
你们在哪里的理解不同？

\exercise{可信度}{个人}{15}

\cite{Fink2013}描述了让老师在学习者眼中可信的三样东西：

\begin{description}

\item[胜任感：]
  对学科的掌握，
  表现为能解释复杂想法、
  或引用他人的工作。

\item[可信赖：]
  把学习者的最大利益放在心上。
  这可以通过给出个性化的反馈、
  对评分决定给出合理的解释、
  以及一视同仁地对待所有学习者来体现。

\item[活力：]
  对学科的热情（\chapref{s:performance}）。

\end{description}

描述你讲课时做的一件归入每个类别的事，
然后描述一件你没做、但应该做的事。

\exercise{衡量有效性}{个人}{15}

\cite{Kirk1994}定义了评估培训的四个层次：

\begin{description}

\item[反应：]
  学习者对培训感觉如何？

\item[学习：]
  他们实际学到了多少？

\item[行为：]
  他们因此改变了多少行为？

\item[结果：]
  那些行为上的改变，如何影响了他们的产出
  或他们所在群体的产出？

\end{description}

为了评估你教什么、怎么教，你在每个层次上都做了什么？
有哪些你能做、却还没做的事？

\exercise{反对意见与反驳}{思考—结对—分享}{15}

你决定不再问学习者你的课有没有用，
因为你知道他们的回答
和他们实际学到多少之间没有相关性（\secref{s:pck-now}）。
于是你提出了四条建议，
每一条都被你的同事驳回了：

\begin{description}

\item[看看他们会不会把课推荐给朋友。]
  这凭什么比问他们对课的感觉更有意义？

\item[最后给他们一次考试。]
  但学习者在一天结束时知道多少，
  并不能很好地预测他们两三个月后还记得多少，
  而且任何形式的期末考试都会让课变得紧张得多。

\item[两三个月后再给他们一次考试。]
  对散养式学习者来说这几乎不可能，
  而且没从工作坊里获得什么的人，
  参与后续活动的可能性也更低，
  所以这样收集到的反馈会有偏。

\item[看看他们会不会继续用学到的东西。]
  在学习者电脑上装间谍软件是招人反感的，
  那这个要怎么实施呢？

\end{description}

先自己想办法，
给出对这些反对意见的回答，
然后把你的回答分享给同伴，
讨论你们各自想出的办法。
做完之后，
把你最看好的办法分享给全班。
